% Attention masks: sliding-window (local) vs. global, and what one token stores in each
\newcommand{\maskgrid}[2]{% x offset, window (0 = global)
  \foreach \q in {1,...,10} {
    \foreach \k in {1,...,10} {
      \pgfmathtruncatemacro{\vis}{(\k <= \q) && ((#2 == 0) || (\q - \k < #2)) ? 1 : 0}
      \ifnum\vis=1
        \fill[black!60] ({#1+(\k-1)*0.28},{-(\q-1)*0.28}) rectangle ++(0.26,-0.26);
      \else
        \draw[black!25] ({#1+(\k-1)*0.28},{-(\q-1)*0.28}) rectangle ++(0.26,-0.26);
      \fi
    }
  }
}
\begin{tikzpicture}[font=\small]
\maskgrid{0}{3}
\maskgrid{4.4}{0}
\node[above] at (1.4,0.1) {local 層（window size $w$）};
\node[above] at (5.8,0.1) {global 層};
\node[rotate=90, font=\scriptsize] at (-0.3,-1.4) {query 位置};
\node[font=\scriptsize] at (1.4,-3.05) {key 位置};
\node[font=\scriptsize] at (5.8,-3.05) {key 位置};
% per-token KV contents
\begin{scope}[xshift=9.0cm, yshift=-0.2cm, font=\scriptsize]
  \node[anchor=west] at (0,0) {\textbf{1 トークン・1 層の KV（31B）}};
  \node[anchor=west] at (0,-0.55) {local:};
  \draw[fill=black!60] (1.3,-0.7) rectangle ++(1.6,0.3); \node at (2.1,-0.55) {\color{white}K};
  \draw[fill=white] (2.9,-0.7) rectangle ++(1.6,0.3); \node at (3.7,-0.55) {V};
  \node[anchor=west] at (1.3,-1.0) {16 head $\times$ 256 dim $\times$ 2};
  \node[anchor=west] at (0,-1.6) {global:};
  \draw[fill=black!60] (1.3,-1.75) rectangle ++(0.8,0.3); \node at (1.7,-1.6) {\color{white}K};
  \draw[fill=white] (2.1,-1.75) rectangle ++(0.8,0.3); \node at (2.5,-1.6) {V};
  \node[anchor=west] at (1.3,-2.05) {4 head $\times$ 512 dim $\times$ 2};
  \node[anchor=west, align=left] at (0,-2.75) {global の V は K と同じprojectionから作る（K=V）が，\\normalizationと RoPE が違うので別々に格納される．\\local は最新 $w$ トークン分，global は全トークン分を保持};
\end{scope}
\end{tikzpicture}
